\documentclass[paper=a4, parskip=half-, ngerman, fontsize=11pt]{scrreprt}
\usepackage[T1]{fontenc}

\usepackage[ngerman]{babel}
\babelprovide[hyphenrules=ngerman-x-latest]{ngerman}

\usepackage{csquotes}
\usepackage[backend=biber]{biblatex}
\addbibresource{Proseminar.bib}

\usepackage{graphicx}
\usepackage{url}
\usepackage{amsmath}
\usepackage{booktabs}
\usepackage{xcolor}

\usepackage{standalone}

% for aligned captions
\usepackage{subcaption}

\usepackage{tikz}
\usetikzlibrary{decorations.markings}
\usepackage[europeanresistors]{circuitikz}
\usetikzlibrary{arrows.meta}
\usetikzlibrary{calc}
\usetikzlibrary{3d}

\usepackage{pgfplots}
\pgfplotsset{compat=1.8}
\usepgfplotslibrary{groupplots}


\begin{document}

% Titelseite
\begin{titlepage}
    \begin{center}
        % Unilogo oben
        \includegraphics[width=0.5\textwidth]{logo_fernuni_hagen.png}\\[2cm]

        {\LARGE \textbf{Übertragung von Signalen auf elektrischen Leitungen}}\\[2cm]

        \textbf{Proseminar Mathematik in der Technik}\\
        \textbf{Modulnummer:} 61711\\
        Fakultät für Mathematik + Informatik\\[0.5cm]

        \begin{tabbing}
            \hspace{6cm} \= \kill
            \textbf{Name:} \> Sven Schmidt \\
            \textbf{Matrikelnummer:} \> 4125169 \\
            \textbf{Abgabedatum:} \> \today \\
            \textbf{Prüfer:} \> PD Dr.-Ing. Stefan Helfert \\
        \end{tabbing}

        \vfill

        {\large FernUniversität in Hagen}\\
        {\large Sommersemester 2025}
    \end{center}
\end{titlepage}

\chapter{Einführung}

\chapter{Modell des Zweidrahtleiters}

In dieser Ausarbeitung beschränken wir uns auf die Untersuchung der Ausbreitung elektromagnetischer Wellen in
Zweidrahtleitern (Paralleldraht), die schematisch in der folgenden Abbildung gezeigt.
\begin{figure}[!h]
    \begin{center}
        \includegraphics[width=0.5\textwidth]{images/Leiter.png}
        \caption{Elektrische Leiter mit ihren Anwendungen. Entnommen aus~\cite{FernuniSkript}.}
        \label{Leiter}
    \end{center}
\end{figure}

Aus den Gesetzen der Elektro- und Magnetostatik ist bekannt, dass sich beim Durchfließen eines elektrischen Stroms um
einen Leiter ein magnetisches Feld ausbildet (Abbildung~\ref{Felder1}). Variiert dieser Strom zeitlich, dann
verändert sich mit ihm das magnetische Feld und dies wiederum bewirkt die Ausbildung eines elektrischen Feldes, welches
sich als induzierte Gegenspannung bemerkbar macht.
\begin{figure}[!ht]
%    \centering
    \begin{subfigure}[b]{0.49\textwidth}
        \centering
        \documentclass{standalone}

\usepackage{tikz}
\usetikzlibrary{decorations.markings}

\usepackage{amsmath}

\begin{document}

\begin{tikzpicture}

\draw[help lines] grid (10,10);

%position of bottom wire
\coordinate (A1) at (5,3);

%position of top wire
\coordinate (A2) at (5,7);

%point defining the top-most mag. field line
\coordinate (start) at (8,10);
\coordinate (through) at (5,6);
\coordinate (end) at (2,10);


% node with negative charge
\draw[fill=orange!10, draw=orange] (A1) circle[radius=0.4] node[anchor=north] {$-i$};
\node[circle, fill=black, inner sep=2pt]  at (A1) {};

% node with positive charge
\draw[fill=orange!10, draw=orange] (A2) circle[radius=0.4] node[anchor=south] {$+i$};
\node at (A2) {\Large$\boldsymbol{\times}$};

% magnetic field lines
\draw[thick, postaction={decorate}, decoration={markings, mark=at position 0.5 with {\arrow{<}}}] (2,8.5) arc (-180:0:3cm and 3cm);
\draw[thick, postaction={decorate}, decoration={markings, mark=at position 0.75 with {\arrow{<}}}] (5,7) circle[radius=0.8];
\draw[thick, postaction={decorate}, decoration={markings, mark=at position 0.5 with {\arrow{>}}}] (8,5) -- (2,5);
\draw[thick, postaction={decorate}, decoration={markings, mark=at position 0.5 with {\arrow{<}}}] (2,1.5) arc (180:0:3cm and 3cm);
\draw[thick, postaction={decorate}, decoration={markings, mark=at position 0.25 with {\arrow{>}}}] (5,3) circle[radius=0.8];


\end{tikzpicture}

\end{document}
        \caption{\color{red}Caption for the first picture}
        \label{Felder1}
    \end{subfigure}%
    \hfill
    \begin{subfigure}[b]{0.49\textwidth}
        \centering
        \documentclass{standalone}
\usepackage[europeanresistors]{circuitikz}
\begin{document}

    \begin{circuitikz}
        % Left and right vertical lines
        \draw (0,0) -- (0,2);
        \draw (3,0) -- (3,2);

        % Top horizontal line with resistor
        \draw (0,1) -- (1,1)
        to[R, l=R] (2,1)
        -- (3,1);

        % Bottom horizontal line with capacitor
        \draw (0,0) -- (1,0)
        to[C, l=C] (2,0)
        -- (3,0);
    \end{circuitikz}

\end{document}

        \caption{\color{red}Caption for the second picture}
        \label{Felder2}
    \end{subfigure}
    \caption{\color{red}Overall caption for the figure}
\end{figure}
Analog geht mit einer an den Leitern anliegende Spannung ein elektrischen Feld einher (Abbildung~\ref{Felder2}).
Ändert sich dieses, so hat dies die Ausbildung eines magnetischen Feldes zur Folge, welches mit dem des Stroms
wechselwirkt.


Im Folgenden nehmen wir an, dass der Leiter parallel verläuft, der Abstand $d$ zwischen den Leitern und der Querschnitt
des Leiters viel kleiner ist, als die Wellenlänge $\lambda$ der auftretenden Signale,
\[ d \ll \lambda = \frac{c_{0}}{f \sqrt{\epsilon_{r}}}. \]
In diesem Fall müssen wir nicht die Maxwellschen Gleichungen
lösen -- statt dessen können wir mit der Methode der Ersatzschaltbilder die Signalübertragung direkt mit den
beteiligten Strömen und Spannungen beschreiben, ohne mit den elektrischen und magnetischen Feldern
argumentieren zu müssen.

Es ist $c_{0}$ die Lichtgeschwindigkeit im Vakuum, $f$ die Frequenz des Signals und $\epsilon_{r}$ die relative
Dielektrizitätskonstante bezeichnet. Wir machen hier die vereinfachende Annahme, dass $\epsilon_{r}$ eine (konstante)
Eigenschaft des Mediums ist, in dem sich das Signal ausbreitet. Im Vakuum gilt $\epsilon_{r} = 1$.


Wir betrachten nun ein kleines, infinitesimales Leitungsstück $\mathrm{d}x$,
\begin{figure}[!h]
    \begin{center}
        \documentclass{standalone}
\usepackage[europeanresistors]{circuitikz}
\begin{document}

    \begin{circuitikz}
        % Left and right vertical lines
        \draw (0,0) -- (0,2);
        \draw (3,0) -- (3,2);

        % Top horizontal line with resistor
        \draw (0,1) -- (1,1)
        to[R, l=R] (2,1)
        -- (3,1);

        % Bottom horizontal line with capacitor
        \draw (0,0) -- (1,0)
        to[C, l=C] (2,0)
        -- (3,0);
    \end{circuitikz}

\end{document}

        \caption{Elektrische Leiter mit ihren Anwendungen. Entnommen aus~\cite{LeitungenUndFilter}.}
        \label{Leitung1}
    \end{center}
\end{figure}

Der Fluss $\Phi$ des magnetischen Feldes $\textbf{H}$ bewirkt eine Induktivität $L_{x}$
\[ \Phi_{x} = i \cdot L_{x} \], wobei $i$ der Strom ist.
Das elektrische Feld $\textbf{E}$ induziert Ladungen auf der Oberfläche des Leiters und damit eine Kapazität $C_{x}$,
die über die Spannung $u$ wie folgt zusammenhängt,
\[ Q_{x} = C_{x} \cdot u . \]
Die gesamte Leitung stellen wir uns aus diesen Leitungsstücken zusammengesetzt vor wie auf Abbildung~\ref{Leitung2}
dargestellt.
\begin{figure}[!h]
    \begin{center}
        \documentclass{standalone}

\usepackage{tikz}
\usepackage[europeanresistors]{circuitikz}
\usetikzlibrary{decorations.markings}
\usetikzlibrary{arrows.meta}
\usetikzlibrary{calc}

\begin{document}

\begin{circuitikz}

\tikzset{
    charge/.style={fill=orange!10, draw=orange, circle, radius=0.4},
    fieldline/.style={postaction={decorate}, decoration={markings}},
    arrow/.style={decoration={mark=at position #1 with {\arrow{Stealth}}}},
    reversearrow/.style={decoration={mark=at position #1 with {\arrowreversed{Stealth}}}}
}

\def\diagram{

%\draw[Stealth-Stealth] (2, 4.25) -- (6, 4.25) node[pos=0.5, circle, fill=white] {d$x$};

\coordinate (A1) at (1,3);
\coordinate (A2) at (3,3);
\coordinate (B1) at (1,0);
\coordinate (B2) at (3,0);
\coordinate (C1) at (5,3);
\coordinate (C2) at (8,3);
\coordinate (D2) at (8,0);
\coordinate (E1) at (10,3);
\coordinate (E2) at (10,0);
\coordinate (E3) at (12,3);
\coordinate (E4) at (12,0);


\draw[{Circle[open, fill=white]}-, fieldline, arrow=0.5] (A1) -- (A2) node[pos=0.5, above] {$i$}
node[pos=0.2] (A3) {};
\draw[{Circle[open, fill=white]}-] (B1) -- (B2) node[pos=0.2] (B3) {};
\draw[-{Circle[open, fill=white]}] (B2) -- (D2);
\draw[-Stealth, shorten <= 1ex, shorten >= 1ex] (A3) -- (B3) node[pos=0.5, right] {$u$};

% resistor
\draw (A2) to[R, l=$R_{x}$, fill=black] (C1);

% inductor
\ctikzset{inductor=american}
\draw (C1) to[L, l=$L_{x}$] (C2);

% capacitor
\draw (C2) to[C, l=$C_{x}$] (D2);

\draw[-{Circle[open, fill=white]}] (C2) -- (E1);
\draw[-{Circle[open, fill=white]}] (D2) -- (E2);

% G
\draw (E1) to[R, l=$G_{x}$] (E2);

\draw[-{Circle[open, fill=white]}] (E1) -- (E3);
\draw[-{Circle[open, fill=white]}] (E2) -- (E4);

}

\diagram

\end{circuitikz}

\end{document}

        \caption{Elektrische Leiter mit ihren Anwendungen. Entnommen aus~\cite{LeitungenUndFilter}.}
        \label{Leitung2}
    \end{center}
\end{figure}
Die führt auf die folgenden qualitativen Eigenschaften: Wir eine Spannung am Leitungsanfang angelegt, muss zunächst die
erste Kapazität $Q_{s_{1}}$ aufgeladen werden. Erst dann bildet sich eine Spannung an der Längsinduktivität $L_{s_{1}}$
an. Der aufgebaute Strom wird nun $Q{s_{2}}$ aufladen etc., bis zum Leiterende. Die Ausbreitung des Einschaltvorgangs
findet mit endlicher Geschwindigkeit statt.

Ein elektrischer Strom, der durch einen Leiter fließt, erfährt einen Widerstand $R_{s}$ auf einer Strecke
$\mathrm{d}s$. Wenn wir annehmen, dass das den Leiter umgebende Dielektrikum leitend ist, dann bezeichnen wir dies mit
dem Querleitwert
$G_{s}$ auf der Strecke $\mathrm{d}s$. Die Abbildung~\ref{Leitung3}
\begin{figure}[!h]
    \begin{center}
        \includegraphics[width=0.5\textwidth]{images/Leitungsstueck.png}
        \caption{Elektrische Leiter mit ihren Anwendungen. Entnommen aus~\cite{LeitungenUndFilter}.}
        \label{Leitung3}
    \end{center}
\end{figure}



\chapter{Die Telegraphenleitungen}
Abbildung~\ref{Leitung4} zeigt nochmals ein infinitesimales Leitungsstück $\mathrm{d}s$ eines verlustbehafteten
Leiters. Am Leitungseingang liegen Strom $i$ und Spannung $u$ an, am Leitungsende $i + \mathrm{d}i$ und $u +
\mathrm{d}u$.
\begin{figure}[!h]
    \begin{center}
        \includegraphics[width=0.5\textwidth]{images/Leitungsstueck2.png}
        \caption{Einschaltbild eines infinitesimalen Leitungsstücks einer verlustbehafteten Leitung. Entnommen
        aus~\cite{LeitungenUndFilter}.}
        \label{Leitung4}
    \end{center}
\end{figure}

Wir definieren auf einem Stück $\mathrm{d}s$

\begin{tabular}{l l}
    den Widerstandsbelag & $R^{\prime} = \frac{R_{s}}{s}$ \\
\addlinespace
    den Induktivitätsbelag & $L^{\prime} = \frac{L_{s}}{s}$ \\
\addlinespace
    den Kapazitätsbelag & $C^{\prime} = \frac{C_{s}}{s}$ \\
\addlinespace
    und den Leitungsbelag & $G^{\prime} = \frac{G_{s}}{s}$. \\
\end{tabular}

Der Widerstand $R^{\prime}$ und die Induktivität $L^{\prime}$ verursachen einen Spannungsabfall von
\[
i \, R^{\prime} \, \mathrm{d}z + \frac{\partial i}{\partial t} \, L^{\prime} \, \mathrm{d} z
\]
entlang von $\mathrm{d}s$. Mit der Kirchhoffschen Maschengleichung {\color{red}REF} folgt für die Spannung
\begin{equation}
u(z, t) = R^{\prime} \, i(z, t) + L^{\prime} \, \frac{\partial i(z, t)}{\partial t} + u(z + \Delta z, t).
\end{equation}
Durch Querleitwert $G_{s}$ und Querkapazität $C_{s}$ fließt der Strom
\[
u \, G^{\prime} \, \mathrm{d}z + \frac{\partial u}{\partial t} \, C^{\prime} \, \mathrm{d} z.
\]
Mit der Kirchhoffschen Knotenregel {\color{red}REF} folgt
\begin{equation}
    i(z, t) = G^{\prime} \, u(z, t) + C^{\prime} \, \frac{\partial u(z + \Delta z, t)}{\partial t} + i(z + \Delta z, t).
\end{equation}
Wir formen nun beide Gleichungen um,
\begin{align}
\frac{u(z + \Delta z, t) - u(z, t)}{\Delta z} &= \frac{R^{\prime}}{\Delta z} \, i(z, t) + \frac{L^{\prime}}{\Delta z}
\,
\frac{\partial i(z, t)}{\partial t} \\
\frac{i(z + \Delta z, t) - i(z, t)}{\Delta z} &= \frac{G^{\prime}}{\Delta z} \, u(z, t) + \frac{C^{\prime}}{\Delta z}
\, \frac{\partial u(z + \Delta z, t)}{\partial t}.
\end{align}
Mit
\[
\lim_{\Delta x \to 0} \frac{u(x+\Delta x, t)-u(x, t)}{\Delta x} = \frac{\partial u(x, t)}{\partial x} \quad \text{und}
\quad \lim_{\Delta x \to 0} \frac{\Delta R}{\Delta x} = R' \quad \text{folgt}
\]
\begin{align}
    \frac{\partial u}{\partial z} &= -\left(R^{\prime} + L^{\prime}\frac{\partial}{\partial t}\right)i \label{eq:Dgl1}
    \\[1ex]
    \frac{\partial i}{\partial z} &= -\left(G^{\prime} + C^{\prime}\frac{\partial}{\partial t}\right)u \label{eq:Dgl2}
\end{align}
Dies ist ein System partieller linearer Differentialgleichungen 1.~Ordnung mit konstanten Koeffizienten. Sie werden
auch Differentialgleichungen der elektrischen Leitung genannt.
Gleichung~\eqref{eq:Dgl1} beschreibt den Spannungsabfall verursacht durch den Leitungswiderstand $R^{\prime}$ und
Induktivität $L^{\prime}$, siehe Abbildung~\ref{Leitung4}.
Die zweite Gleichung~\eqref{eq:Dgl2} beschreibt den Stromfluss zwischen den Leitern, verursacht durch Kapazität
$C^{\prime}$ und Querleitwert $G^{\prime}$.

Oftmals werden diese Gleichungen in die sogenannten Telegraphenleitungen weiter umgeformt. Dazu leiten wir beide
Gleichungen einmal nach $x$ und einmal nach $t$ ab,
\begin{align}
    \frac{\partial^{2} u}{\partial x^2} &= -R^{\prime} \, \frac{\partial i}{\partial x} - L^{\prime} \,
    \frac{\partial^2 i}{\partial t \,
    \partial x} \label{eq:Dgl3} \\[1ex]
    \frac{\partial^{2} i}{\partial x^2} &= -G^{\prime} \, \frac{\partial u}{\partial x} - C^{\prime} \,
    \frac{\partial^2 u}{\partial t \,
    \partial x} \label{eq:Dgl4} \\[1ex]
    \frac{\partial^{2} u}{\partial x \, \partial t} &= -R^{\prime} \, \frac{\partial i}{\partial t} - L^{\prime} \,
    \frac{\partial^2
    i}{\partial t^{2}} \label{eq:Dgl5} \\[1ex]
    \frac{\partial^{2} i}{\partial x \, \partial t} &= -G^{\prime} \, \frac{\partial u}{\partial t} - C^{\prime} \,
    \frac{\partial^2
    u}{\partial t^{2}} \label{eq:Dgl6}
\end{align}
Wir setzen nun \eqref{eq:Dgl2} und \eqref{eq:Dgl6} in \eqref{eq:Dgl3} und \eqref{eq:Dgl1} und \eqref{eq:Dgl5} in
\eqref{eq:Dgl4} ein und erhalten
\begin{align}
    \frac{\partial^{2} u}{\partial x^{2}} &= R^{\prime} G^{\prime} u(x,t) + (R^{\prime} C^{\prime} + L^{\prime}
    G^{\prime}) \frac{\partial u(x, t)}{\partial t} + L^{\prime} C^{\prime} \frac{\partial^{2} u(x,t)}{\partial t^{2}}
     \label{eq:Tele1} \\[1.5ex]
    \frac{\partial^{2} i}{\partial x^{2}} &= R^{\prime} G^{\prime} i(x,t) + (R^{\prime} C^{\prime} + L^{\prime}
    G^{\prime}) \frac{\partial i(x, t)}{\partial t} + L^{\prime} C^{\prime} \frac{\partial^{2} i(x, t)}{\partial t^{2}}
\end{align}
Im Allgemeinen wird man diese partiellen Differentialgleichungen nur numerisch lösen können. Im Folgenden werden wir
einige Spezialfälle betrachten, die es uns erlauben, das qualitative Verhalten der Signalausbreitung zu beschreiben.


\section{Die Gleichungen der verlustlosen Leitung}
\label{VerlustlosesModel}
Da die Telegraphenleitungen i.A. nur nummerisch gelöst werden können, wollen wir hier eine vereinfachende Annahme
machen: wir nehmen an, dass die Leitung verlustlos ist, d.h. wir setzen $R^{\prime} = G^{\prime} = 0$. Dies erlaubt es
uns, die Gleichungen exakt zu lösen und somit das qualitative Verhalten der Signalausbreitung zu beschreiben. Die
Gleichungen \eqref{eq:Dgl1} und \eqref{eq:Dgl2} vereinfachen sich zu
\begin{align}
    \frac{\partial u}{\partial z} &= - L^{\prime} \, \frac{\partial i}{\partial t} \label{eq:Dgl7} \\[1ex]
    \frac{\partial i}{\partial z} &= - C^{\prime} \, \frac{\partial u}{\partial t} \label{eq:Dgl8} .
\end{align}
Jede dieser Gleichungen hat die Form der Wellengleichung. Sie sind homogene partielle Differentialgleichungen, deren
Lösung aus zwei Partiallösungen
\begin{align}
u(x, t) &= u_{v}(v t - x) + u_{r}(v t + x) \label{eq:AllgEq1} \\[1ex]
i(x, t) &= i_{v}(v t - x) + i_{r}(v t + x) \label{eq:AllgEq2}.
\end{align}
Wir nehmen hier stets $v > 0$ an. In diesem Fall handelt es sich bei $u_{v}$ bzw. $i_{v}$ um
vorwärts laufende Wellen, bei $u_{r}$ und $i_{r}$ um rücklaufende Wellen. Wir betrachten dazu
Abbildung~\ref{VorwaertsWelle}.
\begin{figure}[!h]
    \begin{center}
%        \includegraphics[width=0.75\textwidth]{images/VorwaertsWelle.png}
        \documentclass[tikz]{standalone}
\usepackage{pgfplots}
\pgfplotsset{compat=1.8,
    scale only axis,
    xlabel near ticks,
    ylabel near ticks,
    every axis title shift=3pt}

\usetikzlibrary{arrows.meta}
\usepgfplotslibrary{groupplots}

\begin{document}

\begin{tikzpicture}

\pgfplotsset{xmin=0, xmax=15, ymin=0, ymax=2.5,
    axis lines=left,
    xlabel=$x$,
    ylabel=$\textbf{u}$,
    ylabel style={rotate=-90, at={(current axis.above origin)}, anchor=east},
    xlabel style={at={(axis description cs:1,-0.1)},anchor=north},
    xtick={0},
    yticklabels={},
    % scale y
    y=0.75cm
}

% place outside of plots so negative x shift works
\node[anchor=east] at (-0.1cm, 0.5cm) {$t = t_{0}$};
\node[anchor=east] at (-0.1cm, -2.6cm) {$t = t_{1} = t_{0} + \frac{x_{1}}{v}$};


\begin{groupplot}[group style={group size=1 by 2,
horizontal sep=0pt,
vertical sep=1cm}]

\nextgroupplot[]

\addplot[smooth, samples=100, thick, blue, domain=0:15] coordinates {(0,0) (1,2.3) (2,1.75) (3,2) (4,0) (5,0)}
node[pos=0.6, anchor=south west] {\color{black}$u(v t_{0} - x)$} node[pos=0.75] (A1) {};

\draw[-Stealth] (A1) -- +(1cm,0) node[anchor=west] {$v$};


\nextgroupplot[xtick={0, 7}, xticklabels={0, $x_{1}$},]

\addplot[smooth, samples=100, thick, blue] coordinates {(7,0) (8,2.3) (9,1.75) (10,2) (11,0) (12,0)}
node[pos=0.6, anchor=south west] {\color{black}$u(v t_{1} - x)$} node[pos=0.75] (A1) {};

\draw[-Stealth] (A1) -- +(1cm,0) node[anchor=west] {$v$};

\end{groupplot}

\end{tikzpicture}


\end{document}
        \caption{Einschaltbild eines infinitesimalen Leitungsstücks einer verlustbehafteten Leitung. Entnommen
            aus~\cite{LeitungenUndFilter}.}
        \label{VorwaertsWelle}
    \end{center}
\end{figure}
Sei dazu $t_{0} = 0$, $t_{1} > 0$ und $x_{0} > 0$. Dann ist die Welle zum Zeitpunkt $t_{1} > 0$ die Strecke $x = x_{1}
= v t_{1} + x_{0}$ vorangeschritten. Es gilt $u(x_{0}, t_{0}) = u_{v}(-x_{0})$ und
\begin{align*}
    u(x_{1}, t_{1}) &= u_{v}(v t_{1} - x_{1}) \\
                    &= u_{v}(v t_{1} - (v t_{1} + x_{0})) \\
                    &= u_{v}(- x_{0}) \\
                    &= u(x_{0}, t_{0}).
\end{align*}
Wegen $x_{1} > x_{0}$ ist dies also die vorwärts laufende Welle. Analog beschreiben $u_{r}$ bzw. $i_{r}$ jeweils eine
rücklaufende Welle.


\subsubsection{Wellengeschwindigkeit}

Setzen wir die allgemeine Lösung $u_{v}(v t - x)$ in die Gleichungen \eqref{eq:Dgl7} und \eqref{eq:Dgl8} ein, so
erhalten wir mit
\begin{equation*}
    \frac{\partial u_{v}}{\partial x} = \frac{\partial u_{v}}{\partial w} \frac{\partial w}{\partial x} = -1
    \quad \text{und} \quad
    \frac{\partial u_{v}}{\partial t} = \frac{\partial u_{v}}{\partial w} \frac{\partial w}{\partial t} = v
\end{equation*}
\begin{equation*}
    \frac{\partial^{2} u_{v}}{\partial x^{2}} = -1 \frac{\partial}{\partial x} \left( \frac{\partial u_{v}}{\partial w}
    \right) = -1 \frac{\partial^{2} u_{v}}{\partial w^{2}} \frac{\partial w}{\partial x} = \frac{\partial^{2}
    u_{v}}{\partial w^{2}}
\end{equation*}
und
\begin{equation*}
    \frac{\partial^{2} u_{v}}{\partial t^{2}} = v \frac{\partial}{\partial t} \left( \frac{\partial u_{v}}{\partial w}
    \right) = v \frac{\partial^{2} u_{v}}{\partial w^{2}} \frac{\partial w}{\partial t} =v^{2} \frac{\partial^{2}
        u_{v}}{\partial w^{2}}.
\end{equation*}
Es folgt somit
\begin{equation*}
    \frac{\partial^{2} u_{v}}{\partial x^{2}} = L^{\prime} \, C^{\prime} \, \frac{\partial^{2} u_{v}}{\partial t^{2}}
    \Leftrightarrow
    \frac{\partial^{2} u_{v}}{\partial w^{2}} = L^{\prime} \, C^{\prime} \, v^{2} \frac{\partial^{2} u_{v}}{\partial
    w^{2}}.
\end{equation*}
Nachdem wir die partielle Ableitung herauskürzen, erhalten wir also
\begin{equation}
    v = \frac{1}{\sqrt{L^{\prime} \, C^{\prime}}} \label{eq:WellenGeschw1}.
\end{equation}
überraschend ist nun, dass die Ausbreitungsgeschwindigkeit des Signals im verlustlosen Leiter durch die Parameter
$L^{\prime}$ und $C^{\prime}$ abhängt (und damit vom Isolatormaterial zwischen den Leitern), jedoch nicht vom
Leitermaterial selbst!


\subsubsection{Wellenwiderstand}

Wir kommen nun zu einer weiteren wichtigen Größe, dem Wellenwiderstand.
Setzen wir Gleichungen \eqref{eq:AllgEq1}, \eqref{eq:AllgEq2} und \eqref{eq:WellenGeschw1} in \eqref{eq:Dgl7} und
\eqref{eq:Dgl8} ein, so folgt
\begin{align*}
    \frac{\partial u}{\partial x} &= \frac{\partial u_{v}}{\partial w_{-}} \frac{\partial w_{-}}{\partial x}
    + \frac{\partial u_{r}}{\partial w_{+}} \frac{\partial w_{+}}{\partial x}
    =
    \frac{\partial u_{v}}{\partial w_{-}} \left(- \frac{1}{v} \right) + \frac{\partial u_{r}}{\partial w_{+}}
    \frac{1}{v} \\[1ex]
    \frac{\partial i}{\partial t} &= \frac{\partial i_{v}}{\partial w_{-}} - \frac{\partial i_{r}}{\partial w_{+}}
\end{align*}
und
\begin{align*}
    \frac{\partial i}{\partial x} &= \frac{\partial i_{v}}{\partial w_{-}} \frac{\partial w_{-}}{\partial x}
    + \frac{\partial i_{r}}{\partial w_{+}} \frac{\partial w_{+}}{\partial x}
    =
    \frac{\partial i_{v}}{\partial w_{-}} \left(- \frac{1}{v} \right) - \frac{\partial i_{r}}{\partial w_{+}}
    \frac{1}{v} \\[1ex]
    \frac{\partial u}{\partial t} &= \frac{\partial u_{v}}{\partial w_{-}}
    \underbrace{\frac{\partial w_{-}}{\partial t} }_{=1}
    -
    \frac{\partial u_{r}}{\partial w_{+}} \underbrace{\frac{\partial w_{+}}{\partial t} }_{=1}
\end{align*}
und somit
\begin{alignat*}{3}
    \frac{\partial u}{\partial x} &= \frac{1}{v} \left( \frac{\partial u_{r}}{\partial w_{+}} -
    \frac{\partial u_{v}}{\partial w_{-}} \right) &= - L^{\prime} \left( \frac{\partial i_{v}}{\partial w_{-}} -
    \frac{\partial i_{r}}{\partial w_{+}} \right) &= L^{\prime} \, \frac{\partial i}{\partial t} \\[1ex]
    \frac{\partial i}{\partial x} &= - \frac{1}{v} \left( \frac{\partial i_{r}}{\partial w_{+}} +
    \frac{\partial i_{v}}{\partial w_{-}} \right) &= - C^{\prime} \left( \frac{\partial u_{v}}{\partial w_{+}} +
    \frac{\partial u_{r}}{\partial w_{-}} \right) &= C^{\prime} \, \frac{\partial u}{\partial t}.
\end{alignat*}
Da die Funktionen $u_{v}$ und $u_{r}$ bzw. $i_{v}$ und $i_{r}$ unabhängig sind, folgt
\begin{equation}
    - \frac{1}{v} \, \frac{\text{d} u_{v}}{\text{d} w_{-}} = - L^{\prime} \, \frac{\text{d} i_{v}}{\text{d} w_{-}}.
\end{equation}
Integrieren beider Seiten liefert
\begin{equation}
    u_{v}(w_{-}) = v \, L^{\prime} \, i_{v}(w_{-}) + c_{v} = \sqrt{\frac{L^{\prime}}{C^{\prime}}} \, i_{v}(w_{-}) +
    c_{v},
\end{equation}
und
\begin{equation}
    u_{r}(w_{+}) = v \, L^{\prime} \, i_{r}(w_{+}) + c_{r} = \sqrt{\frac{L^{\prime}}{C^{\prime}}} \, i_{r}(w_{+}) +
    c_{r},
\end{equation}
wobei $c_{v}$ und $c_{r}$ Konstanten sind. Setzen wir $c_{v} = c_{r} = 0$, dann erhalten wir
\begin{equation}
    \frac{u_{v}(x, t)}{i_{v}(x, t)} = \frac{u_{r}(x, t)}{i_{r}(x, t)} = \sqrt{\frac{L^{\prime}}{C^{\prime}}} = Z_{w}.
\end{equation}
Diese Konstante hat die Maßeinheit eines Widerstandes, wir nennen $Z_{w}$ folglich den Wellenwiderstand der Leitung.
Kommen wir nun zur verlustbehafteten Leitung.

\section{Verlustbehaftetes Model}
Wir wollen jetzt für einen weiteren Spezialfall die Telegraphenleitungen lösen. Wir nehmen an, dass das System im
eingeschwungenen Zustand ist. Wir nehmen weiter an, dass die auftretenden Signale aus einer einzigen sin-förmigen
Schwingung fester Frequenz bestehen. I.A. Fall zerlegt man die auftretenden Signale in ihre einzelnen
Frequenzkomponenten und betrachtet diese zunächst getrennt.

Wir stellen Strom und Spannung mittels sogenannter Phasoren $\underline{I}$ und $\underline{U}$ dar, um hervorzuheben,
dass es sich um komplexe Größen handelt:
\begin{align*}
    i(x, t) &= \mathrm{Re} \left( \underline{I}(x) \exp(j \omega t) \right) \\
    u(x, t) &= \mathrm{Re} \left( \underline{U}(x) \exp(j \omega t) \right)
\end{align*}
mit $\omega = 2 \pi f$. Es handelt sich hierbei um stationäre sin-förmige Signale, bei denen Ort- und Zeitabhängigkeit
entkoppelt sind.

Das Einsetzen der Phasoren in die Gleichungen \eqref{eq:Dgl1} und \eqref{eq:Dgl2} liefert
\begin{align}
    \frac{\text{d} \underline{U}}{\text{d} x} &= - \left( R^{\prime} + j \omega L^{\prime} \right)
    \underline{I} \label{eq:ResDgl1} \\[1ex]
    \frac{\text{d} \underline{I}}{\text{d} x} &= - \left( G^{\prime} + j \omega C^{\prime} \right) \underline{U}
\end{align}
Wir eliminieren $\underline{I}$ aus \eqref{eq:ResDgl1} durch nochmaliges Ableiten der Gleichung nach $x$:
\begin{align}
    \frac{\text{d}^{2} \underline{U}}{\text{d} x^{2}} &= \left( R^{\prime} + j \omega L^{\prime} \right)
    \left( G^{\prime} + j \omega C^{\prime} \right) \underline{U} \label{eq:ResDgl3}.
\end{align}
Verglichen mit \eqref{eq:Tele1} entfällt in \eqref{eq:ResDgl3} die Zeitabhängigkeit. Damit haben wir eine gewöhnliche
Differentialgleichung 2.~Ordnung mit konstanten Koeffizienten, wir wir sie in Abschnitt~\ref{VerlustlosesModel} bereits
begegnet sind.

Für den Phasor $\underline{I}$ erhalten wir mit ähnlichen Argumenten
\begin{align}
    \frac{\text{d}^{2} \underline{I}}{\text{d} x^{2}} &= \left( R^{\prime} + j \omega L^{\prime} \right)
    \left( G^{\prime} + j \omega C^{\prime} \right) \underline{I}.
\end{align}
Dies sind die Wellengleichungen der verlustbehafteten Leitung. Sie werden in der Regel wie folgt geschrieben:
\begin{align}
    \frac{\text{d}^{2} \underline{U}}{\text{d} x^{2}} &= \gamma^{2} \underline{U} \\[1ex]
    \frac{\text{d}^{2} \underline{I}}{\text{d} x^{2}} &= \gamma^{2} \underline{I} \\[1ex]
    \gamma^{2} &= \left( R^{\prime} + j \omega L^{\prime} \right) \left( G^{\prime} + j \omega C^{\prime} \right).
\end{align}
Man nennt $\gamma$ die Wellenausbreitungsgeschwindigkeit. Verglichen zum verlustlosen Model ist dies eine komplexe
Größe, was dazu führt, dass im verlustbehafteten Fall ein Dämpfungsterm in der Lösung auftritt.

Die allgemeine Lösungen haben die Form
\[ \underline{U}_{1} \exp(- \gamma x) \quad \text{und} \quad \underline{U}_{2} \exp(\gamma x), \] wobei
$\underline{U}_{1}$ und $\underline{U}_{2}$ Integrationskonstanten sind, die durch die Anfangsbedingungen festgelegt
werden.
Es gilt somit für die Spannung
\begin{equation}
    \underline{U}(x) = \underline{U}_{1} \exp(- \gamma x) + \underline{U}_{2} \exp(\gamma x) \label{eq:VerlustU}.
\end{equation}
Die Gleichung für den Strom erhalten wir, wenn wir \eqref{eq:VerlustU} in \eqref{eq:ResDgl1} einsetzen,
\begin{equation}
    \underline{I}(x) = \frac{-1}{R^{\prime} + j \omega L^{\prime}} \frac{\text{d} \underline{U}}{\text{d} x}
    = \frac{1}{Z}
    \left(
    \vphantom{\sqrt{1}{Z}}
    \underline{U}_{1} \exp(- \gamma x) + \underline{U}_{2} \exp(\gamma x) \right),
\end{equation}
wobei wir
\begin{equation}
    Z = \sqrt{\frac{G^{\prime} + j \omega C^{\prime}}{R^{\prime} + j \omega L^{\prime}}}
\end{equation}
eingeführt haben. $Z$ hat die Dimension eines Widerstandes und wird daher Wellenwiderstand oder Wellenimpedanz
genannt.

Die Konstanten $\underline{U}_{1}$ und $\underline{U}_{2}$ werden in der Regel durch die angelegte Spannung am
Leiteranfang festgelegt. Gilt z.B.
\begin{align*}
    \underline{U}(0) &= \underline{U}_{a} = \underline{U}_{1} + \underline{U}_{2} \\[1ex]
    \underline{I}(0) &= \underline{I}_{a} = \frac{1}{Z}
    \left(
    \vphantom{\sqrt{1}{Z}}
    \underline{U}_{1} - \underline{U}_{2} \right),
\end{align*}
dann folgt für die Integrationskonstanten
\[ \underline{U}_{1} = \frac{\underline{U}_{a} + Z \underline{I}_{a}}{2} \quad \text{und} \quad \underline{U}_{2} =
\frac{\underline{U}_{a} - Z \underline{I}_{a}}{2}. \]
Damit lauten die Gleichungen für Spannung und Strom entlang der Leitung
\begin{align}
    \underline{U}(x) &=
    \frac{1}{2} \left( \underline{U}_{a} + Z \underline{I}_{a} \right) \exp (- \gamma x)
    +
    \frac{1}{2} \left( \underline{U}_{a} - Z \underline{I}_{a} \right) \exp (\gamma x) \\[1ex]
    \underline{I}(x) &=
    \frac{1}{2} \left( \frac{\underline{U}_{a}}{Z} + \underline{I}_{a} \right) \exp (- \gamma x)
    -
    \frac{1}{2} \left( \frac{\underline{U}_{a}}{Z} - \underline{I}_{a} \right) \exp (\gamma x).
\end{align}



\section{Abschlüsse}


%\bibliographystyle{plain} % oder ein anderer Stil, z.B. alpha, abbrv, unsrt
%\bibliography{Proseminar} % Name Ihrer .bib-Datei OHNE Endung .bib
\printbibliography

\end{document}